% title: 

\begin{document}
\section*{2026/8/7}

现在是8月7日凌晨三点三十六分，我再一次失眠了。最近的事情非常非常多，包括但不限于19号国赛的答辩(但是18号就得到哈尔滨)，16号之前提交常微分方程的作业，20号之前学完密码分析学，开学之前写完复变函数论的作业，还有准备密码技术竞赛的作品。

这么多事情......已经同疯掉无异了。


\section*{2026/8/4}

我本来没打算在这个暑假折腾我的个人网站，本来的计划是完成网安数学站的建立，然后在开学的时候申请一下在各个年级的年级大会上展示一下，以此来维护和延续网安数学站这个项目。

网安数学站的本意是为了通俗易懂地教会学院内的数学课程，并且联系相关的一些数学书籍来获得更好的理解，但是因为我始终坚持自己手搓文档而不用AI来生成，再加上学院内的数学课程实在是太多了，导致我的进度非常非常缓慢，甚至是有点停滞不前，搞得我自己有些道心崩溃，再加上这个暑假的安排也逐渐变多，再过个十几天还要去参加一个比赛的答辩，一个半月后要提交另外一个比赛的比赛作品，确实是没有时间来更新网安数学站的内容了。

而且我认真思考了一下，做网安数学站的好处其实只有一个，就是这个网站对所有人都是有益的，我十分肯定这一点，从不同课本中筛选出最好的理解方式再加上通俗易懂的语言，对所有人它都是一个很好的平台，但是做网安数学站的坏处实在是太多了，首先是我个人，写这些文档要花掉难以想象的时间和精力，我之前也差不多统计过，累计写二十多个小时完不成课本上一章的内容，太耗费时间精力了，再加上，这个项目做出来也只是第一步，困难的是后面的维护工作，学院数学课程是不断更新的，我不可能毕业了之后还要维护这个网站，做不到也不可能，而且我觉得很难找到后面的人来接手这个项目，维护起来太累而且收益很低，所以我只能决定暂时放弃这个项目。

以后可以每个周末抽一点零碎时间来完成文档的更新，然后缓慢坚持下去，这样才有可能把这个网站给做出来，也只是有可能。



\end{document}
